\documentclass[tikz,border=0pt]{standalone}
%% =====================================================================
%%  gridkit.tex — shared TikZ utilities for GridKit documentation figures
%% ---------------------------------------------------------------------
%%  Usage (from docs/Figures/<Area>/<Name>/diagram.tex):
%%
%%      \documentclass[tikz,border=12pt]{standalone}
%%      %% =====================================================================
%%  gridkit.tex — shared TikZ utilities for GridKit documentation figures
%% ---------------------------------------------------------------------
%%  Usage (from docs/Figures/<Area>/<Name>/diagram.tex):
%%
%%      \documentclass[tikz,border=12pt]{standalone}
%%      %% =====================================================================
%%  gridkit.tex — shared TikZ utilities for GridKit documentation figures
%% ---------------------------------------------------------------------
%%  Usage (from docs/Figures/<Area>/<Name>/diagram.tex):
%%
%%      \documentclass[tikz,border=12pt]{standalone}
%%      \input{../../utils/gridkit.tex}
%%      \begin{document}
%%      \begin{tikzpicture}[gk]
%%          ...
%%      \end{tikzpicture}
%%      \end{document}
%%
%%  Naming: colors gk<Name>, styles gk<name>, macros \gk<Name>.
%%  `\documentclass[tikz,...]{standalone}` already loads TikZ; this file
%%  only adds packages, libraries, colors, styles, and macros.
%%
%%  Extending this file:
%%    * New measurement   -> add a \newcommand{\gk<Name>} in MACROS, keeping
%%                           the signature {<coords...>}{<label>}.
%%    * New color / role   -> add one \definecolor/\colorlet in PALETTE and
%%                           expose it through a gk<name> style.
%%    * New recurring shape-> add a gk<name>/.style or a macro.
%%    * Never redefine an existing public name (keep figures back-compatible).
%% =====================================================================

% ---- Packages & libraries -------------------------------------------
\usepackage{amsmath}
\usetikzlibrary{arrows.meta, positioning, calc, decorations.pathreplacing}

% ---- Palette --------------------------------------------------------
\definecolor{gkConductor}{RGB}{20,70,130}   % energized conductor (blue)
\definecolor{gkSteel}{RGB}{95,95,100}       % steel: poles, ground, dimensions
\colorlet{gkImage}{gkConductor!45}          % mirror-image conductor (muted)
\colorlet{gkGuide}{gkSteel!70}              % chords / reference guide lines
\colorlet{gkEarth}{gkSteel!12}              % earth fill below the ground line

% ---- Picture-wide defaults (activate with [gk]) ---------------------
\tikzset{
  gk/.style = { >={Stealth[length=2.4mm]}, line cap=round },
}

% ---- Named styles ---------------------------------------------------
\tikzset{
  % Block-diagram family (matches existing EMT diagrams)
  gkline/.style    = {semithick, rounded corners=6pt},
  gkblock/.style   = {draw, semithick, fill=white,
                      minimum height=1cm, minimum width=2.3cm},
  gksum/.style     = {draw, semithick, fill=white, circle, inner sep=2.6pt},
  gksignal/.style  = {circle, fill, inner sep=1.4pt, label distance=3pt},
  gksig/.style     = {font=\small},
  % Physical-geometry family
  gkcond/.style    = {gkConductor, line width=1.7pt},
  gkconddot/.style = {gkConductor, fill=gkConductor},
  gkimage/.style   = {gkImage, line width=1.2pt, dashed},
  gkimagedot/.style= {gkImage, fill=gkImage},
  gkguide/.style   = {gkGuide, dashed},
  gkdim/.style     = {<->, thin, gkSteel},
  gkdimlbl/.style  = {fill=white, inner sep=1.5pt},
}

% ---- Drawing macros -------------------------------------------------
% Earth baseline with hatch ticks:  \gkground{xLeft}{xRight}{y}
\newcommand{\gkground}[3]{%
  \draw[gkSteel, thick] (#1,#3) -- (#2,#3);%
  \pgfmathsetmacro{\gkstep}{#1+0.6}%
  \foreach \gkx in {#1,\gkstep,...,#2}{%
    \draw[gkSteel, thin] (\gkx,#3) -- (\gkx-0.28,#3-0.32);}%
}
% Steel support pole + top attachment dot:  \gkpole{x}{yBottom}{yTop}
\newcommand{\gkpole}[3]{%
  \draw[gkSteel, line width=2.2pt] (#1,#2) -- (#1,#3);%
  \fill[gkSteel] (#1,#3) circle (2.4pt);%
}
% Horizontal dimension line:  \gkdimh{x1}{x2}{y}{label}
\newcommand{\gkdimh}[4]{%
  \draw[gkdim] (#1,#3) -- (#2,#3) node[midway, gkdimlbl]{#4};%
}
% Vertical dimension line:  \gkdimv{x}{y1}{y2}{label}
\newcommand{\gkdimv}[4]{%
  \draw[gkdim] (#1,#2) -- (#1,#3) node[midway, right=2.5pt]{#4};%
}
% Styled conductor point marker:  \gkdot{coord}  (cartesian or polar, e.g. 78:5)
\newcommand{\gkdot}[1]{\fill[gkconddot] (#1) circle (3.2pt);}
% Radial dimension from a centre:  \gkradius{cx}{cy}{angle}{len}{label}
\newcommand{\gkradius}[5]{%
  \draw[->, thin, gkSteel] (#1,#2) -- ($ (#1,#2) + (#3:#4) $)
       node[pos=0.58, gkdimlbl]{#5};%
}
% Angle arc between two directions:  \gkangle{cx}{cy}{angA}{angB}{radius}{label}
\newcommand{\gkangle}[6]{%
  \draw[thin, gkSteel] ($ (#1,#2) + (#3:#5) $) arc (#3:#4:#5);%
  \pgfmathsetmacro{\gkmid}{(#3+#4)/2}%
  \node[gksig] at ($ (#1,#2) + (\gkmid:#5+0.5) $) {#6};%
}
% Filled conductor disk:  \gkdisk{x}{y}{r}
\newcommand{\gkdisk}[3]{\fill[gkConductor] (#1,#2) circle (#3);}
% Hollow conductor annulus:  \gkannulus{x}{y}{rOuter}{rInner}
\newcommand{\gkannulus}[4]{%
  \fill[gkConductor] (#1,#2) circle (#3);%
  \fill[white] (#1,#2) circle (#4);%
  \draw[gkConductor, semithick] (#1,#2) circle (#3) (#1,#2) circle (#4);%
}
% Catenary sag arc (midspan at y=0):  \gkcatenary{xL}{xR}{a}
%   plots z = a[cosh(x/a) - 1]; supports sit at z = a[cosh(|x|/a)-1]
\newcommand{\gkcatenary}[3]{%
  \draw[gkcond] plot[domain=#1:#2, samples=180] (\x,{#3*(cosh(\x/#3)-1)});%
}

%%      \begin{document}
%%      \begin{tikzpicture}[gk]
%%          ...
%%      \end{tikzpicture}
%%      \end{document}
%%
%%  Naming: colors gk<Name>, styles gk<name>, macros \gk<Name>.
%%  `\documentclass[tikz,...]{standalone}` already loads TikZ; this file
%%  only adds packages, libraries, colors, styles, and macros.
%%
%%  Extending this file:
%%    * New measurement   -> add a \newcommand{\gk<Name>} in MACROS, keeping
%%                           the signature {<coords...>}{<label>}.
%%    * New color / role   -> add one \definecolor/\colorlet in PALETTE and
%%                           expose it through a gk<name> style.
%%    * New recurring shape-> add a gk<name>/.style or a macro.
%%    * Never redefine an existing public name (keep figures back-compatible).
%% =====================================================================

% ---- Packages & libraries -------------------------------------------
\usepackage{amsmath}
\usetikzlibrary{arrows.meta, positioning, calc, decorations.pathreplacing}

% ---- Palette --------------------------------------------------------
\definecolor{gkConductor}{RGB}{20,70,130}   % energized conductor (blue)
\definecolor{gkSteel}{RGB}{95,95,100}       % steel: poles, ground, dimensions
\colorlet{gkImage}{gkConductor!45}          % mirror-image conductor (muted)
\colorlet{gkGuide}{gkSteel!70}              % chords / reference guide lines
\colorlet{gkEarth}{gkSteel!12}              % earth fill below the ground line

% ---- Picture-wide defaults (activate with [gk]) ---------------------
\tikzset{
  gk/.style = { >={Stealth[length=2.4mm]}, line cap=round },
}

% ---- Named styles ---------------------------------------------------
\tikzset{
  % Block-diagram family (matches existing EMT diagrams)
  gkline/.style    = {semithick, rounded corners=6pt},
  gkblock/.style   = {draw, semithick, fill=white,
                      minimum height=1cm, minimum width=2.3cm},
  gksum/.style     = {draw, semithick, fill=white, circle, inner sep=2.6pt},
  gksignal/.style  = {circle, fill, inner sep=1.4pt, label distance=3pt},
  gksig/.style     = {font=\small},
  % Physical-geometry family
  gkcond/.style    = {gkConductor, line width=1.7pt},
  gkconddot/.style = {gkConductor, fill=gkConductor},
  gkimage/.style   = {gkImage, line width=1.2pt, dashed},
  gkimagedot/.style= {gkImage, fill=gkImage},
  gkguide/.style   = {gkGuide, dashed},
  gkdim/.style     = {<->, thin, gkSteel},
  gkdimlbl/.style  = {fill=white, inner sep=1.5pt},
}

% ---- Drawing macros -------------------------------------------------
% Earth baseline with hatch ticks:  \gkground{xLeft}{xRight}{y}
\newcommand{\gkground}[3]{%
  \draw[gkSteel, thick] (#1,#3) -- (#2,#3);%
  \pgfmathsetmacro{\gkstep}{#1+0.6}%
  \foreach \gkx in {#1,\gkstep,...,#2}{%
    \draw[gkSteel, thin] (\gkx,#3) -- (\gkx-0.28,#3-0.32);}%
}
% Steel support pole + top attachment dot:  \gkpole{x}{yBottom}{yTop}
\newcommand{\gkpole}[3]{%
  \draw[gkSteel, line width=2.2pt] (#1,#2) -- (#1,#3);%
  \fill[gkSteel] (#1,#3) circle (2.4pt);%
}
% Horizontal dimension line:  \gkdimh{x1}{x2}{y}{label}
\newcommand{\gkdimh}[4]{%
  \draw[gkdim] (#1,#3) -- (#2,#3) node[midway, gkdimlbl]{#4};%
}
% Vertical dimension line:  \gkdimv{x}{y1}{y2}{label}
\newcommand{\gkdimv}[4]{%
  \draw[gkdim] (#1,#2) -- (#1,#3) node[midway, right=2.5pt]{#4};%
}
% Styled conductor point marker:  \gkdot{coord}  (cartesian or polar, e.g. 78:5)
\newcommand{\gkdot}[1]{\fill[gkconddot] (#1) circle (3.2pt);}
% Radial dimension from a centre:  \gkradius{cx}{cy}{angle}{len}{label}
\newcommand{\gkradius}[5]{%
  \draw[->, thin, gkSteel] (#1,#2) -- ($ (#1,#2) + (#3:#4) $)
       node[pos=0.58, gkdimlbl]{#5};%
}
% Angle arc between two directions:  \gkangle{cx}{cy}{angA}{angB}{radius}{label}
\newcommand{\gkangle}[6]{%
  \draw[thin, gkSteel] ($ (#1,#2) + (#3:#5) $) arc (#3:#4:#5);%
  \pgfmathsetmacro{\gkmid}{(#3+#4)/2}%
  \node[gksig] at ($ (#1,#2) + (\gkmid:#5+0.5) $) {#6};%
}
% Filled conductor disk:  \gkdisk{x}{y}{r}
\newcommand{\gkdisk}[3]{\fill[gkConductor] (#1,#2) circle (#3);}
% Hollow conductor annulus:  \gkannulus{x}{y}{rOuter}{rInner}
\newcommand{\gkannulus}[4]{%
  \fill[gkConductor] (#1,#2) circle (#3);%
  \fill[white] (#1,#2) circle (#4);%
  \draw[gkConductor, semithick] (#1,#2) circle (#3) (#1,#2) circle (#4);%
}
% Catenary sag arc (midspan at y=0):  \gkcatenary{xL}{xR}{a}
%   plots z = a[cosh(x/a) - 1]; supports sit at z = a[cosh(|x|/a)-1]
\newcommand{\gkcatenary}[3]{%
  \draw[gkcond] plot[domain=#1:#2, samples=180] (\x,{#3*(cosh(\x/#3)-1)});%
}

%%      \begin{document}
%%      \begin{tikzpicture}[gk]
%%          ...
%%      \end{tikzpicture}
%%      \end{document}
%%
%%  Naming: colors gk<Name>, styles gk<name>, macros \gk<Name>.
%%  `\documentclass[tikz,...]{standalone}` already loads TikZ; this file
%%  only adds packages, libraries, colors, styles, and macros.
%%
%%  Extending this file:
%%    * New measurement   -> add a \newcommand{\gk<Name>} in MACROS, keeping
%%                           the signature {<coords...>}{<label>}.
%%    * New color / role   -> add one \definecolor/\colorlet in PALETTE and
%%                           expose it through a gk<name> style.
%%    * New recurring shape-> add a gk<name>/.style or a macro.
%%    * Never redefine an existing public name (keep figures back-compatible).
%% =====================================================================

% ---- Packages & libraries -------------------------------------------
\usepackage{amsmath}
\usetikzlibrary{arrows.meta, positioning, calc, decorations.pathreplacing}

% ---- Palette --------------------------------------------------------
\definecolor{gkConductor}{RGB}{20,70,130}   % energized conductor (blue)
\definecolor{gkSteel}{RGB}{95,95,100}       % steel: poles, ground, dimensions
\colorlet{gkImage}{gkConductor!45}          % mirror-image conductor (muted)
\colorlet{gkGuide}{gkSteel!70}              % chords / reference guide lines
\colorlet{gkEarth}{gkSteel!12}              % earth fill below the ground line

% ---- Picture-wide defaults (activate with [gk]) ---------------------
\tikzset{
  gk/.style = { >={Stealth[length=2.4mm]}, line cap=round },
}

% ---- Named styles ---------------------------------------------------
\tikzset{
  % Block-diagram family (matches existing EMT diagrams)
  gkline/.style    = {semithick, rounded corners=6pt},
  gkblock/.style   = {draw, semithick, fill=white,
                      minimum height=1cm, minimum width=2.3cm},
  gksum/.style     = {draw, semithick, fill=white, circle, inner sep=2.6pt},
  gksignal/.style  = {circle, fill, inner sep=1.4pt, label distance=3pt},
  gksig/.style     = {font=\small},
  % Physical-geometry family
  gkcond/.style    = {gkConductor, line width=1.7pt},
  gkconddot/.style = {gkConductor, fill=gkConductor},
  gkimage/.style   = {gkImage, line width=1.2pt, dashed},
  gkimagedot/.style= {gkImage, fill=gkImage},
  gkguide/.style   = {gkGuide, dashed},
  gkdim/.style     = {<->, thin, gkSteel},
  gkdimlbl/.style  = {fill=white, inner sep=1.5pt},
}

% ---- Drawing macros -------------------------------------------------
% Earth baseline with hatch ticks:  \gkground{xLeft}{xRight}{y}
\newcommand{\gkground}[3]{%
  \draw[gkSteel, thick] (#1,#3) -- (#2,#3);%
  \pgfmathsetmacro{\gkstep}{#1+0.6}%
  \foreach \gkx in {#1,\gkstep,...,#2}{%
    \draw[gkSteel, thin] (\gkx,#3) -- (\gkx-0.28,#3-0.32);}%
}
% Steel support pole + top attachment dot:  \gkpole{x}{yBottom}{yTop}
\newcommand{\gkpole}[3]{%
  \draw[gkSteel, line width=2.2pt] (#1,#2) -- (#1,#3);%
  \fill[gkSteel] (#1,#3) circle (2.4pt);%
}
% Horizontal dimension line:  \gkdimh{x1}{x2}{y}{label}
\newcommand{\gkdimh}[4]{%
  \draw[gkdim] (#1,#3) -- (#2,#3) node[midway, gkdimlbl]{#4};%
}
% Vertical dimension line:  \gkdimv{x}{y1}{y2}{label}
\newcommand{\gkdimv}[4]{%
  \draw[gkdim] (#1,#2) -- (#1,#3) node[midway, right=2.5pt]{#4};%
}
% Styled conductor point marker:  \gkdot{coord}  (cartesian or polar, e.g. 78:5)
\newcommand{\gkdot}[1]{\fill[gkconddot] (#1) circle (3.2pt);}
% Radial dimension from a centre:  \gkradius{cx}{cy}{angle}{len}{label}
\newcommand{\gkradius}[5]{%
  \draw[->, thin, gkSteel] (#1,#2) -- ($ (#1,#2) + (#3:#4) $)
       node[pos=0.58, gkdimlbl]{#5};%
}
% Angle arc between two directions:  \gkangle{cx}{cy}{angA}{angB}{radius}{label}
\newcommand{\gkangle}[6]{%
  \draw[thin, gkSteel] ($ (#1,#2) + (#3:#5) $) arc (#3:#4:#5);%
  \pgfmathsetmacro{\gkmid}{(#3+#4)/2}%
  \node[gksig] at ($ (#1,#2) + (\gkmid:#5+0.5) $) {#6};%
}
% Filled conductor disk:  \gkdisk{x}{y}{r}
\newcommand{\gkdisk}[3]{\fill[gkConductor] (#1,#2) circle (#3);}
% Hollow conductor annulus:  \gkannulus{x}{y}{rOuter}{rInner}
\newcommand{\gkannulus}[4]{%
  \fill[gkConductor] (#1,#2) circle (#3);%
  \fill[white] (#1,#2) circle (#4);%
  \draw[gkConductor, semithick] (#1,#2) circle (#3) (#1,#2) circle (#4);%
}
% Catenary sag arc (midspan at y=0):  \gkcatenary{xL}{xR}{a}
%   plots z = a[cosh(x/a) - 1]; supports sit at z = a[cosh(|x|/a)-1]
\newcommand{\gkcatenary}[3]{%
  \draw[gkcond] plot[domain=#1:#2, samples=180] (\x,{#3*(cosh(\x/#3)-1)});%
}


% =====================================================================
%  Path model: an ordered GIS route P_1,...,P_{Np} on the earth sphere.
%  Each interval P_n -> P_{n+1} contributes the great-circle (haversine)
%  arc length s_n = R_oplus c_n; the route is a smooth chain of geodesic
%  segments.  The globe is an offset, cropped, soft-lit sphere.
%
%  Orthographic projection; each segment is sampled by spherical linear
%  interpolation (slerp) of its endpoint unit vectors.
% =====================================================================

% modern globe palette (cool blue-gray, lit from upper-left)
\definecolor{gkGlobeLo}{RGB}{237,243,248}   % highlight (near-white)
\definecolor{gkGlobeHi}{RGB}{151,174,196}   % body / terminator
\tikzset{gkgrid/.style = {white, line width=0.4pt, opacity=0.5}}

\def\Rs{5.2}     % sphere screen radius (large)
\def\ox{2.7}     % globe centre x (offset right)
\def\oy{-1.3}    % globe centre y (offset down)
\def\vlat{16}    % view latitude  (deg)
\def\vlon{20}    % view longitude (deg)
\def\Nseg{20}    % samples per geodesic

\begin{document}

% ---- orthographic view basis (east, north, out) at (vlat,vlon) -----
\pgfmathsetmacro{\eastx}{-sin(\vlon)}
\pgfmathsetmacro{\easty}{ cos(\vlon)}
\pgfmathsetmacro{\eastz}{0}
\pgfmathsetmacro{\northx}{-sin(\vlat)*cos(\vlon)}
\pgfmathsetmacro{\northy}{-sin(\vlat)*sin(\vlon)}
\pgfmathsetmacro{\northz}{ cos(\vlat)}
\pgfmathsetmacro{\outx}{cos(\vlat)*cos(\vlon)}
\pgfmathsetmacro{\outy}{cos(\vlat)*sin(\vlon)}
\pgfmathsetmacro{\outz}{sin(\vlat)}

% project (lat,lon) -> screen (\sx,\sy); \sz>0 means front hemisphere
\newcommand{\proj}[2]{%
  \pgfmathsetmacro{\vX}{cos(#1)*cos(#2)}%
  \pgfmathsetmacro{\vY}{cos(#1)*sin(#2)}%
  \pgfmathsetmacro{\vZ}{sin(#1)}%
  \pgfmathsetmacro{\sx}{\Rs*(\vX*\eastx+\vY*\easty+\vZ*\eastz)}%
  \pgfmathsetmacro{\sy}{\Rs*(\vX*\northx+\vY*\northy+\vZ*\northz)}%
  \pgfmathsetmacro{\sz}{\vX*\outx+\vY*\outy+\vZ*\outz}%
}
\newcommand{\waypoint}[3]{\proj{#2}{#3}\coordinate (#1) at (\sx,\sy);}

% great-circle polyline between (lat1,lon1) and (lat2,lon2)
\newcommand{\geoseg}[5]{% #1 lat1 #2 lon1 #3 lat2 #4 lon2 #5 style
  \pgfmathsetmacro{\Ax}{cos(#1)*cos(#2)}%
  \pgfmathsetmacro{\Ay}{cos(#1)*sin(#2)}%
  \pgfmathsetmacro{\Az}{sin(#1)}%
  \pgfmathsetmacro{\Bx}{cos(#3)*cos(#4)}%
  \pgfmathsetmacro{\By}{cos(#3)*sin(#4)}%
  \pgfmathsetmacro{\Bz}{sin(#3)}%
  \pgfmathsetmacro{\dab}{min(1,max(-1,\Ax*\Bx+\Ay*\By+\Az*\Bz))}%
  \pgfmathsetmacro{\Om}{acos(\dab)}%
  \foreach \i in {0,...,\Nseg}{%
    \pgfmathsetmacro{\tt}{\i/\Nseg}%
    \pgfmathsetmacro{\wa}{sin((1-\tt)*\Om)/sin(\Om)}%
    \pgfmathsetmacro{\wb}{sin(\tt*\Om)/sin(\Om)}%
    \pgfmathsetmacro{\vX}{\wa*\Ax+\wb*\Bx}%
    \pgfmathsetmacro{\vY}{\wa*\Ay+\wb*\By}%
    \pgfmathsetmacro{\vZ}{\wa*\Az+\wb*\Bz}%
    \pgfmathsetmacro{\px}{\Rs*(\vX*\eastx+\vY*\easty+\vZ*\eastz)}%
    \pgfmathsetmacro{\py}{\Rs*(\vX*\northx+\vY*\northy+\vZ*\northz)}%
    \coordinate (gc\i) at (\px,\py);%
  }%
  \draw[#5] (gc0) \foreach \i in {1,...,\Nseg}{ -- (gc\i)};%
}

% front-hemisphere parallel (lat = #1) and meridian (lon = #1)
\newcommand{\parallelarc}[1]{%
  \foreach \lo in {-180,-165,...,165}{%
    \pgfmathsetmacro{\loB}{\lo+15}%
    \proj{#1}{\lo}\edef\xa{\sx}\edef\ya{\sy}\edef\za{\sz}%
    \proj{#1}{\loB}\edef\xb{\sx}\edef\yb{\sy}\edef\zb{\sz}%
    \ifdim\za pt>0pt\ifdim\zb pt>0pt
      \draw[gkgrid] (\xa,\ya) -- (\xb,\yb);\fi\fi}}
\newcommand{\meridianarc}[1]{%
  \foreach \la in {-75,-60,...,75}{%
    \pgfmathsetmacro{\laB}{\la+15}%
    \proj{\la}{#1}\edef\xa{\sx}\edef\ya{\sy}\edef\za{\sz}%
    \proj{\laB}{#1}\edef\xb{\sx}\edef\yb{\sy}\edef\zb{\sz}%
    \ifdim\za pt>0pt\ifdim\zb pt>0pt
      \draw[gkgrid] (\xa,\ya) -- (\xb,\yb);\fi\fi}}

\begin{tikzpicture}[gk]

% fixed figure frame (defines crop + bounding box)
\useasboundingbox (-4.6,-3.7) rectangle (4.6,3.7);

\begin{scope}
\clip (-4.6,-3.7) rectangle (4.6,3.7);
\begin{scope}[shift={(\ox,\oy)}]

  % ---- soft-lit sphere: body + diffuse highlight (upper-left) -------
  \pgfmathsetmacro{\hlx}{-0.50*\Rs}
  \pgfmathsetmacro{\hly}{ 0.52*\Rs}
  \pgfmathsetmacro{\hlr}{ 1.35*\Rs}
  \begin{scope}
    \clip (0,0) circle (\Rs);
    \fill[gkGlobeHi] (0,0) circle (\Rs);
    \shade[inner color=gkGlobeLo, outer color=gkGlobeHi] (\hlx,\hly) circle (\hlr);
  \end{scope}

  % ---- graticule (front hemisphere) and crisp limb -----------------
  \foreach \L in {-60,-30,0,30,60}{\parallelarc{\L}}
  \foreach \L in {-60,-30,0,30,60,90}{\meridianarc{\L}}
  \draw[gkGlobeHi!60!black, thin] (0,0) circle (\Rs);

  % ---- smooth GIS route P_1 ... P_8 on the upper-left surface -------
  \waypoint{P1}{36}{-24}
  \waypoint{P2}{40}{-19}
  \waypoint{P3}{42}{-14}
  \waypoint{P4}{40}{ -9}
  \waypoint{P5}{36}{ -4}
  \waypoint{P6}{34}{  1}
  \waypoint{P7}{37}{  5}
  \waypoint{P8}{41}{ 10}

  \geoseg{36}{-24}{40}{-19}{gkcond}
  \geoseg{40}{-19}{42}{-14}{gkcond}
  \geoseg{42}{-14}{40}{ -9}{gkcond}
  \geoseg{40}{ -9}{36}{ -4}{gkcond}
  \geoseg{36}{ -4}{34}{  1}{gkcond}
  \geoseg{34}{  1}{37}{  5}{gkcond}
  \geoseg{37}{  5}{41}{ 10}{gkcond}

  \foreach \p in {P1,P2,P3,P4,P5,P6,P7,P8}{%
    \fill[white] (\p) circle (3.4pt);
    \fill[gkconddot] (\p) circle (2.6pt);}
  \node[gksig, gkConductor, below=3pt]  at (P1) {$P_1$};
  \node[gksig, gkConductor, right=3pt]  at (P8) {$P_{N_p}$};

\end{scope}
\end{scope}

% ---- title / coordinate note (top-left, off the globe) -------------
\node[gksig, gkConductor, anchor=north west] at (-4.4,3.5)
     {$P_n=(\Phi_n,\Lambda_n)$};

% ---- detail callout (lower-left): one segment's arc geometry -------
\begin{scope}[shift={(-2.55,-1.95)}]
  \def\ri{1.4}
  \fill[white] (-1.55,-0.72) rectangle (1.55,2.18);
  \draw[gkSteel, thin, rounded corners=3pt] (-1.55,-0.72) rectangle (1.55,2.18);
  \draw[gkSteel, thin] (0,0) -- (110:\ri);
  \draw[gkSteel, thin] (0,0) -- (70:\ri) node[pos=0.56, right=2pt, gksig, gkSteel]{$R_\oplus$};
  \draw[gkcond, line width=1.6pt] (70:\ri) arc (70:110:\ri);
  \node[gksig] at (90:\ri+0.28) {$s_n$};
  \draw[thin, gkSteel] (70:0.55) arc (70:110:0.55);
  \node[gksig, gkSteel] at (90:0.88) {$c_n$};
  \fill[gkconddot] (110:\ri) circle (2.4pt) node[gkConductor, above left=-2pt and -2pt, gksig]{$P_n$};
  \fill[gkconddot] ( 70:\ri) circle (2.4pt) node[gkConductor, above right=-2pt and -2pt, gksig]{$P_{n+1}$};
  \fill[gkSteel] (0,0) circle (1.5pt) node[below=1.5pt, gksig]{$O$};
\end{scope}

\end{tikzpicture}
\end{document}
